\chapter{Charaktergrundlagen}

Dieses Kapitel beschreibt die Grundlagen eines Charakters, seine Entwicklung
über Level sowie den Umgang mit Talenten und Fähigkeiten.

\section{Charaktererschaffung}

Jeder Charakter beginnt das Spiel auf \textbf{Level 1} und erhält:

\begin{itemize}
	\item 1 Klasse
	\item \textbf{15 \termref[Talentpunkte]{talentpunkte}}
	\item \textbf{15 \termref[Fähigkeitspunkte]{faehigkeitspunkte}}
	\item bis zu \textbf{2 Eigenheiten}
\end{itemize}

Talentpunkte und Fähigkeitspunkte sind voneinander unabhängig
und werden getrennt verwaltet.

\section{Charakterstufe}

Die Charakterstufe (Level) beschreibt den allgemeinen Entwicklungsstand eines Charakters. Das Level des Charakters wird durch die Anzahl der gesamten Fähigkeitspunkte definiert. Punkte, die zum Erreichen eines \textbf{\termref[legendären Rangs (Rang 6)]{legendaerer-rang}} ausgegeben werden, zählen \textbf{nicht} für den Levelaufstieg.

Mit steigenden Leveln erhält ein Charakter:
\begin{itemize}
	\item zusätzliche \termref[Fähigkeitspunkte]{faehigkeitspunkte}
	\item Bonuspunkte auf bestimmten Stufen
	\item mehr Aktions- oder Manapunkte
\end{itemize}

Ab \textbf{Level 2} kann zusätzlich zur Hauptklasse eine \textbf{beliebige Nebenklasse}
gewählt werden.

\begin{center}
	\begin{tabular}{lccccc}
		\toprule
		\textbf{Level} & 1 & 2 & 3 & 4 & 5 \\
		\midrule
		Fähigkeitspunkte (gesamt) & 10 & 20 & 35 & 45 & 60 \\
		Bonusfähigkeitspunkte & -- & +5 & -- & +5 & -- \\
		Aktions-/Manapunkte & 50 & 60 & 70 & 80 & 90 \\
		\bottomrule
	\end{tabular}
\end{center}

Die zusätzlichen Bonusfähigkeitspunkte auf Level 2 und 4 werden \textbf{zusätzlich}
zu allen anderen Punkten vergeben.



\section{Klassen und Fähigkeiten}

Klassen und \termref[Fähigkeiten]{faehigkeiten} (Skills und Zauber) müssen jeweils \textbf{einzeln erlernt und gesteigert}
werden.

Dabei gilt:
\begin{itemize}
	\item \termref[Magierklassen]{magierklassen} lernen Zauber
	\item \termref[Nichtmagische Klassen]{nicht-magische-klassen} lernen Skills
	\item Basisfähigkeiten können von allen Klassen erlernt werden
	\item Die Hauptklasse kann ab Level 1 gewählt werden
	\item Eine Nebenklasse kann ab Level 2 gewählt werden
\end{itemize}

\section{Talentpunkte und Benchmark}

Talentpunkte dienen der Steigerung von Talenten und deren Rängen.

Zur Orientierung der Charakterstärke dient folgende Benchmark-Tabelle:

\begin{center}
	\begin{tabular}{lcccc}
		\toprule
		\textbf{Benchmark} & Normal & Beruf & Elite & Max \\
		\midrule
		Talentpunkte & 11 & 22 & 33 & 44 \\
		\bottomrule
	\end{tabular}
\end{center}

\subsection*{Talente als Grundlage des Charakters}

\termref[Talente]{talente} bilden das grundlegende Kompetenzgerüst eines Charakters.
Sie bestimmen, wie gut ein Charakter in zentralen Bereichen wie Kampf,
Bewegung, Wissen oder sozialer Interaktion ist.

Im nächsten Kapitel werden Talente, ihre Ränge, Kosten sowie legendären
Effekte im Detail beschrieben.

\section{Ressourcen}

\term{Ressourcen}{ressourcen} beschreiben die körperlichen, geistigen und besonderen Reserven
eines Charakters. Sie bestimmen, wie lange ein Charakter handeln, kämpfen,
zaubern oder besondere Effekte nutzen kann.

\subsection*{Lebenspunkte (LP)}

\begin{rulebox}{Lebenspunkte}
	\term{Lebenspunkte}{lebenspunkte-lp} repräsentieren die körperliche Belastbarkeit und
	Vitalität eines Charakters.
\end{rulebox}

\textbf{Berechnung:}
\[
\text{LP} = 20 + \text{Stärke-Wurfzahl}
\]

\textbf{Regeln:}
\begin{itemize}
	\item Wird ein Charakter auf \textbf{0 LP} reduziert, \textbf{stirbt} er.
	\item Schaden wird erst auf Schutzressourcen angewendet
	(siehe Kampfregeln), bevor Lebenspunkte verloren gehen.
\end{itemize}

\textbf{Regeneration:}  
Pro vollständiger Ruhephase (z.\,B. Nachtruhe) regeneriert ein Charakter
Lebenspunkte in Höhe von:
\[
2 \times \text{Stärke-Rang}
\]

\textbf{Heilung:}  
Tränke, Zauber und bestimmte Skills können Lebenspunkte sofort oder über Zeit
wiederherstellen.

\subsection*{Aktionspunkte (AP) und Manapunkte (MP)}

\begin{rulebox}{Aktions- und Manapunkte}
	\term{Aktionspunkte}{aktionspunkte-ap} und \term{Manapunkte}{manapunkte-mp} beschreiben die
	handlungs- und zauberbezogenen Reserven eines Charakters.
\end{rulebox}

\begin{itemize}
	\item Auf Level 1 stehen insgesamt \textbf{50 Punkte} zur Verfügung.
	\item Diese können frei auf Aktionspunkte (AP) und Manapunkte (MP)
	aufgeteilt werden.
	\item Mit steigenden Leveln erhöht sich dieser Pool
	(siehe Level-Tabelle).
\end{itemize}

\textbf{Regeneration:}  
Pro vollständiger Ruhephase (z.\,B. Nachtruhe) regeneriert ein Charakter
\textbf{2W6 AP} und \textbf{MP}.

\subsection*{Materialien}

\begin{rulebox}{Materialien}
	\term{Materialien}{ressource-materialien} sind eine abstrakte Ressource,
	die vorbereitete, improvisierte oder verbrauchende Handlungen abbildet.
	Sie stehen für Komponenten, Hilfsmittel, Verbrauchsstoffe oder situative
	Vorbereitungen, ohne als einzelne Gegenstände geführt zu werden.
\end{rulebox}

\textbf{Allgemeine Regeln:}
\begin{itemize}
	\item Materialien sind zu behandeln wie ganz normale Gegenstände. Sie können gefunden, gekauft oder gestohlen werden.
	\item Materialien werden durch \textbf{Fähigkeiten,
		Klassenregeln oder Effekte} verbraucht.
	\item Startmenge von Materialien sind
	\textbf{klassenabhängig} und werden im jeweiligen Klassenkapitel festgelegt. 
\end{itemize}

\textbf{Verwendung:}
\begin{itemize}
	\item Der Verbrauch von Materialien ist in der jeweiligen Fähigkeit
	explizit angegeben.
	\item Hat ein Charakter nicht genügend Materialien, kann die Fähigkeit
	nicht eingesetzt werden.
\end{itemize}

\subsection*{Marken}

\begin{rulebox}{Marken}
	\term{Marken}{marken} sind eine seltene Ressource, die legendäre Effekte
	aktivieren oder verstärken kann.
\end{rulebox}

\begin{itemize}
	\item Marken werden durch das Erreichen einer \textbf{Wurfzahl $\ge 10$}
	bei Proben erhalten.
	\item Ein Charakter kann maximal \textbf{2 Marken} gleichzeitig besitzen.
	\item Marken werden ausschließlich für legendäre Effekte verwendet.
\end{itemize}


\section{Fortschritt durch Abenteuer}

\subsection*{Skill- und Zauberpunkte}

Empfohlen ist der Erhalt von \textbf{3–5 Fähigkeitspunkte}
pro Abenteuer, abhängig von Kampagnenlänge und Spielstil.

Zusätzliche Punkte können durch besondere Leistungen vergeben werden
(z.\,B. Questbelohnungen, gutes Rollenspiel oder kreative Lösungen).

\subsection*{Talentpunkte}

Empfohlen ist der Erhalt von \textbf{2–3 Talentpunkten} pro Abenteuer,
abhängig von Kampagnenlänge und gewünschter Progression.

\section{Legendärer Rang (Rang 6)}

\begin{rulebox}{Legendärer Rang}
	Rang 6 ist der \term{legendären Rang}{legendaerer-rang} einer Disziplin.
	Er verwendet den \textbf{W20} und unterliegt der Ass-Regel wie gewohnt.
	Zudem schaltet Rang 6 für Talente und Fähigkeiten \term{legendäre Effekte}{legendaere-efffekte} frei.
\end{rulebox}

Der Zugang zu Rang 6 ist streng begrenzt und erfordert besondere Voraussetzungen.

\subsection*{Legendär-Slots}

Ein Charakter verfügt über eine stark begrenzte Anzahl legendärer Slots:

\begin{itemize}
	\item \textbf{Talente:} Genau \textbf{ein} Talent darf Rang 6 erreichen.
	\item \textbf{Fähigkeiten (Skills/Zauber):} Genau \textbf{drei} Legendär-Slots:
	\begin{itemize}
		\item 1× Basis
		\item 1× Hauptklasse
		\item 1× Nebenklasse
	\end{itemize}
\end{itemize}

Wird keine Nebenklasse gespielt, kann der entsprechende Slot alternativ
für eine weitere Basis-Fähigkeit verwendet werden.

\subsection*{Voraussetzungen}

Um eine Disziplin auf Rang 6 zu steigern, müssen alle folgenden Bedingungen erfüllt sein:

\begin{enumerate}
	\item Die gewählte Disziplin (Talent oder Fähigkeit) befindet sich auf \textbf{Rang 5}.
	\item Ein passender \textbf{Lehrmeister oder Fokusort} wurde gefunden.
	\item Es wurden mindestens \textbf{drei ausreichende Trainingsphasen} investiert.
\end{enumerate}

\subsection*{Meisterprobe}

Nach Abschluss der Trainingsphasen erfolgt genau \textbf{eine Meisterprobe}
mit dem aktuellen Rang-5-Würfel der Disziplin.

\begin{itemize}
	\item \textbf{Erfolg bei 10+}: Aufstieg zu Rang 6 (W20) und Bindung des Slots.
	\item \textbf{Misserfolg}: Kein Malus. Ein erneuter Versuch ist frühestens
	nach einer weiteren Trainingsphase und einer neuen Szene möglich.
\end{itemize}

Es wird immer nur \textbf{ein einziger Wurf} durchgeführt.

\subsection*{Zugriff auf legendäre Effekte}

Legendäre Effekte dürfen \textbf{ausschließlich} genutzt werden, wenn:
\begin{itemize}
	\item die jeweilige Disziplin selbst Rang 6 erreicht hat und
	\item ein legendärer Slot auf genau diese Disziplin gebunden ist
\end{itemize}

Legendäre Effekte gelten nur für den konkret gewählten Eintrag.
Andere Talente oder Fähigkeiten derselben Klasse profitieren nicht.

Ohne Rang 6 besteht \textbf{kein Zugriff} auf legendäre Effekte.

\subsection*{Lehrmeister und Fokusorte}

Wer einen legendären Rang lehren kann, hängt von der Disziplin ab:

\begin{itemize}
	\item \textbf{Haupt- und Nebenklasse:} Orden, Meister oder Zirkel der jeweiligen Klasse
	\item \textbf{Basis-Fähigkeiten:} Gilden, heilige Stätten oder freie Meister
	\item \textbf{Talente:} Lehrmeister oder Fokusorte mit klarer Talent-Expertise
\end{itemize}

% =========================
% Eigenheiten
% =========================

\section{Eigenheiten}

\term{Eigenheiten}{eigenheiten} sind dauerhafte \textbf{Trade-Offs}, die einen Charakter
mechanisch und erzählerisch prägen.  
Jede Eigenheit besteht immer aus \textbf{einem Vorteil} und \textbf{einem Nachteil}.

Eigenheiten können:
\begin{itemize}
	\item bei der Charaktererschaffung gewählt werden
	\item im Verlauf der Kampagne \textbf{erlangt, verändert oder verloren} werden
\end{itemize}

Eigenheiten sind kein reiner Bonus und kein reiner Malus – sie verändern aktiv,
\textbf{wie} ein Charakter gespielt wird.

\medskip

\begin{rulebox}{Grundregeln zu Eigenheiten}
	\begin{itemize}
		\item Ein Charakter beginnt mit \textbf{1–2 Eigenheiten}.
		\item Eigenheiten haben \textbf{keinen Rang} und skalieren nicht.
		\item Effekte von Eigenheiten sind \textbf{immer aktiv}, sofern nichts anderes
		angegeben ist.
	\end{itemize}
\end{rulebox}

\medskip

\begin{rulebox}{Eigenheiten im Spielverlauf}
	Eigenheiten können sich im Laufe der Kampagne verändern:
	\begin{itemize}
		\item Neue Eigenheiten können durch
		\textbf{dramatische Ereignisse, Entscheidungen oder Konsequenzen}
		entstehen.
		\item Eigenheiten können \textbf{abgeschwächt, verstärkt oder vollständig verloren}
		gehen, wenn sich der Charakter entsprechend entwickelt.
		\item Änderungen an Eigenheiten erfolgen immer in Absprache mit der Spielleitung.
	\end{itemize}
\end{rulebox}

\medskip

\begin{rulebox}{Hinweis zur Balance}
	Eigenheiten sind bewusst \textbf{stark und spürbar}.
	
	Ein Charakter mit vielen Eigenheiten ist nicht automatisch stärker,
	sondern komplexer und konfliktanfälliger zu spielen.
\end{rulebox}



